
%%%%%%%%%%%%%%%%%%%%%%%%%%%%%%%%%%%%%%%%%%%%%%%%%%%%%%%%%%%%%%%%%%%%%%%%%%%%%%%%%%%%%%%%%%%%%%%%%%%%%%%%%%%%%%%%%%%%%%%%%%%%%%%%%%%%%%%%%%%%%%%%%%%%%%%%%%%%%%%%%%%%%%%%%%%%%%%%%%%%%%%%%%%%%%%%%%%%%%%%%%%%%%%%%%%%%%%%%%%%%%%%%%%%%%%%%%%%%%%%%%%%%%%%%%%%%%%%%%%%%%%%%%%%%%%%%%%%%%%%%%%%%%%%%%%%%%%%%%%%%%%%%%%%%%%%%%%%%%%%%%%%



%\begin{comment} % TIPOGRAFICA  
%%%%%%%%%%%% Para ver si lo puedes eliminar porque queda demasidado feo con estas caracteristicas %%%%%%%%%
% Tipo de letra: Arial.  
\usepackage{helvet}  % Usa Helvetica (similar a Arial)
\renewcommand{\familydefault}{\sfdefault}  % Configura la fuente sans-serif por defecto

% Tamaño de letra de título de capítulo: 14 puntos. 
\usepackage{titlesec}
\titleformat{\chapter}{ \centering \fontsize{14}{16}\bfseries}{\Roman{chapter}.}{1em}{}

\titleformat{\section}{\fontsize{11}{13}\bfseries}{\makebox[0.3cm][r]{\thesection}}{1em}{}

% Configurar sangría de subapartados (Tercer nivel, Números arábigos con sangría 0.8 cm)
\titleformat{\subsection}
  {\bfseries}{\makebox[0.8cm][r]{\thesubsection}~}{0pt}{}





%Tamaño de letra de texto dentro de tablas y figuras: 9 puntos.
\captionsetup[figure]{font={small,stretch=1.5}}  % small es 9pt
\captionsetup[table]{font={small,stretch=1.5}}
%\end{comment}

% Tablas y figuras: 1.50 cm. o 0.59 pulgadas. 
\captionsetup{font={stretch=1.5}}

%\setlength{\floatsep}{1.5cm}

\renewcommand{\thechapter}{\Roman{chapter}.}
\renewcommand{\thesection}{\Alph{section}.}
\renewcommand{\thesubsection}{\arabic{subsection}} % Subapartados en números arábigos (1, 2, 3...)
%\renewcommand{\thesection}{\arabic{chapter}.\arabic{section}}


\renewcommand{\footnotesize}{\fontsize{9pt}{10pt}\selectfont}  % Tamaño de fuente en notas
%\end{comment}



%\begin{comment} % INTERLINEADO
% El texto principal es a doble espacio: 2.00 cm o 0.78 pulgadas. 

\doublespacing  % Aplica doble espaciado en todo el documento

% Ecuaciones, desde 1.00 cm. hasta 3.00 cm. dependiendo de la  diagramación
\setlength{\abovedisplayskip}{1.5cm}  % Espacio antes de ecuaciones
\setlength{\belowdisplayskip}{1.5cm}  % Espacio después de ecuaciones
% \setlength{\jot}{1cm}  % Espaciado entre ecuaciones en align



% Notas de páginas: 1.00 cm. o 0.39 pulgadas. 
\setlength{\skip\footins}{1cm}  % Ajusta la distancia entre el texto y las notas al pie


%%%%%%%%%%%%%%%%%%%%%%%%%%%%%%%%%%%%%%%%%%%%%%%%%%%%%%%%%%%%%%%%%%%%%%%%%%%%%%%%%%%%%%%%%%%%%%%%%%%%%%%%%%%%%%%%%%%%%%%%%%%%%%%%%%%%%%%%%%%%%%%%%%%%%%%%%%%%%%%%%%%%%%%%%%%%%%
% numeración arabiga continua
 \counterwithout{equation}{chapter}
 \renewcommand{\theequation}{\arabic{equation}} % Asegura numeración en arábigos
\setcounter{equation}{0} % Inicia la numeración desde 1



% Evitar que las ecuaciones se reinicien en cada capítulo


% Evitar que la numeración se reinicie por capítulo
\counterwithout{table}{chapter}
\counterwithout{figure}{chapter}


%%%%%% Numeración de tablas en números romanos %%
\renewcommand{\thetable}{\Roman{table}} % Usa números romanos para tablas
% Configurar el formato de las tablas para IEEE
\captionsetup[table]{labelfont=normalfont, labelsep=newline}
%%%%%% Numeración de figuras en números romanos %%
\captionsetup[figure]{labelsep=period} % Para cambiar ":" por "."
\renewcommand{\thefigure}{\Roman{figure}}



%%%%%%%%%% Separación entre numero y texto en la listas %%%%%%%%%%%%%
\setlength{\cftfignumwidth}{7.9em}  % Ajusta la anchura del número
\setlength{\cfttabnumwidth}{6.5em}  % Ajusta la separación en la lista de tablas
% Cambiar numeración a romanos SOLO en \listoffigures y \listoftables

% Corregir la numeración en la lista de figuras (usar \cftfigaftersnum)
\renewcommand{\cftfigpresnum}{FIGURA\;} 
\renewcommand{\cftfigaftersnum}{\Roman{figure}.}

\renewcommand{\cfttabpresnum}{TABLA\;}
\renewcommand{\cfttabaftersnum}{\Roman{table}:}



% Sangría en la tabla de contenido
\setlength{\cftchapnumwidth}{0.85cm}
\setlength{\cftsecnumwidth}{0.6cm}     % Espacio reservado para el número de la sección
\setlength{\cftsecindent}{0.3cm} 
\setlength{\cftsubsecindent}{0.8cm} % Sangría de 40pt en subsecciones


\usepackage{titlesec}
\titleformat{\section}
  {\bfseries\itshape}
  {\thesection}
  {1em}
  {}
\titleformat{\subsection}
  {\bfseries\itshape}
  {\thesubsection}
  {1em}
  {}